\documentclass[aspectratio=169,xcolor=dvipsnames]{beamer}
\usetheme{Simple}
\usepackage[T1]{fontenc}
\usepackage[utf8]{inputenc}
\usepackage[english]{babel}
\usepackage{graphicx}
\usepackage{booktabs}
\usepackage{array}
\usepackage{ragged2e}
\newcolumntype{L}[1]{>{\RaggedRight\arraybackslash}p{#1}}
\usepackage{xcolor}
\usepackage{hyperref}

\definecolor{CyberRed}{HTML}{B42332}
\definecolor{CyberBurgundy}{HTML}{7F1D1D}
\definecolor{CyberLight}{HTML}{FCE8EA}
\definecolor{CyberText}{HTML}{292525}
\setbeamercolor{normal text}{fg=CyberText,bg=white}
\setbeamercolor{structure}{fg=CyberRed}
\setbeamercolor{title}{fg=white}
\setbeamercolor{subtitle}{fg=white}
\setbeamercolor{frametitle}{fg=CyberBurgundy,bg=white}
\setbeamercolor{section in toc}{fg=CyberBurgundy}
\setbeamercolor{item}{fg=CyberRed}
\setbeamercolor{subitem}{fg=CyberRed}
\setbeamercolor{enumerate item}{fg=CyberRed}
\setbeamercolor{alerted text}{fg=CyberRed}
\setbeamercolor{block title}{fg=white,bg=CyberBurgundy}
\setbeamercolor{block body}{fg=CyberText,bg=CyberLight}
\setbeamercolor{block title alerted}{fg=white,bg=CyberRed}
\setbeamercolor{block body alerted}{fg=CyberText,bg=CyberLight}
\setbeamercolor{block title example}{fg=white,bg=CyberBurgundy}
\setbeamercolor{block body example}{fg=CyberText,bg=CyberLight}
\setbeamercolor{palette primary}{fg=white,bg=CyberRed}
\setbeamercolor{palette secondary}{fg=white,bg=CyberBurgundy}
\setbeamercolor{palette tertiary}{fg=white,bg=CyberBurgundy}
\setbeamercolor{footline}{fg=CyberBurgundy,bg=white}
\setbeamercolor{navigation symbols}{fg=CyberRed}
\hypersetup{colorlinks=true,urlcolor=CyberRed,linkcolor=CyberBurgundy,citecolor=CyberRed}

% EDIT THE FIELDS BELOW.
\title[Azure AI -- Module 1]{Microsoft Azure AI -- Module 1}
\subtitle{Microsoft Azure -- Practical Approach \\ Azure Virtual Machines as an IaaS solution}
\author[Student / team]{Filip Kłos, Szymon Kaczmarek, Krzysztof Szczepaniak, Jakub Rybczyński, Dorian Koehler}
\institute[SUT]{Faculty of Applied Mathematics\\Silesian University of Technology}
\date{\today}

\begin{document}
	\begin{frame}
		\titlepage
	\end{frame}
	\begin{frame}{Presentation outline}
		\tableofcontents
	\end{frame}
	\section{Introduction}
	\begin{frame}{Selected Azure topic and objectives}
		\begin{itemize}
			\item \textbf{Selected topic:} Azure Virtual Machines (IaaS)
			\item \textbf{Problem context:} A team needs a Linux environment where it can install custom runtimes, configure the OS, and run software that is not easily packaged for a PaaS platform.
			\item \textbf{Why Azure VMs:} They provide the highest degree of control over the VM and its operating system compared with other Azure compute options
			\item \textbf{What will be presented:} VM concepts, portal configuration choices, a practical deployment/inspection, observed results, and cost/security trade-offs.
			\item \textbf{Note:} Azure VMs free you from buying physical hardware, but they do not free you from OS maintenance, patching and software installation remain your responsibility.
		\end{itemize}
	\end{frame}
	\section{Cloud concepts}
	\begin{frame}{Cloud foundations and terminology}
		\begin{itemize}
			\item \textbf{IaaS (Infrastructure as a service)} - You rent compute, manage the OS, runtime, and applications.
			\item \textbf{PaaS (Platform as a service)} - You rent service, don't manage underlying OS
			\item \textbf{Where VMs fit} - Azure VMs are one of several on-demand, scalable compute resource types in Azure
			\item \textbf{Prerequisites} - An Azure subscription, a resource group, a region, permission to create resources.
			\item \textbf{Billing boundary} - The subscription defines billing and access; the resource group collects related resources.
		\end{itemize}
	\end{frame}
	\section{Azure exploration}
	\begin{frame}{Selected service: overview}
		\begin{itemize}
			\item \textbf{Purpose of Azure VMs} - On-demand, scalable compute with virtualization, without purchasing physical hardware.
			\item \textbf{Common use cases} - Development and testing, cloud applications with variable demand, and extending an on-premises network into Azure.
			\item \textbf{Key configuration choices}
			\begin{itemize}
				\item Subscription
				\item Region and resource group
				\item VM name
				\item VM image
				\item VM architecture
				\item VM size/family
				\item Authentication
				\item Inbound port rules
			\end{itemize}
			\item \textbf{Where to find it} Azure portal → Create a resource → search Virtual machine.
		\end{itemize}
	\end{frame}
	\begin{frame}{Azure VM families}
		\textbf{The first cost/performance decision}
		
		\vspace{0.2cm}
		Azure groups VM sizes into \textbf{families} by CPU-to-memory ratio and
		specialised hardware. The wrong family causes overspending
		or underperformance
		
		\vspace{0.3cm}
		\textbf{The name encodes the specification}
		
		\vspace{0.1cm}
		{\footnotesize
			\begin{center}
				\begin{tabular}{@{}ll@{}}
					\toprule
					\textbf{Part} & \textbf{Meaning (example)} \\
					\midrule
					\textbf{D}  & Family: General purpose \\
					\textbf{C}  & Subfamiily \\
					\textbf{a}  & CPU info \\
					\textbf{d}  & Storage \\
					\textbf{s}  & Premium storage \\
					\textbf{v5} & 5th hardware generation \\
					\bottomrule
				\end{tabular}
			\end{center}
		}
	\end{frame}
	\begin{frame}[shrink=10]{Core VM families}
		\centering
		\renewcommand{\arraystretch}{1.4}
		\setlength{\tabcolsep}{4pt}
		\small
		
		\begin{tabular}{@{} L{1.6cm} L{2.2cm} L{2.6cm} L{3.8cm} L{3.2cm} @{}}
			\toprule
			\textbf{Family} & \textbf{Focus} & \textbf{Approx.\ profile} & \textbf{Best for} & \textbf{Cost note} \\
			\midrule
			\textbf{B-series} & Burstable
			& Baseline + CPU credits
			& Dev/test, low-traffic servers, occasional spikes
			& Lowest for idle-heavy use; throttled when credits end \\
			\addlinespace
			\textbf{D-series} & General purpose
			& Balanced ($\sim$4\,GB/vCPU)
			& Web/app servers, microservices
			& Default baseline; steady performance \\
			\addlinespace
			\textbf{E-series} & Memory optimised
			& High RAM ($\sim$8\,GB/vCPU)
			& In-memory DBs, caches, analytics
			& Higher cost/vCPU; pay for RAM \\
			\addlinespace
			\textbf{F-series} & Compute optimised
			& High CPU ($\sim$2\,GB/vCPU)
			& Batch, CI/CD, CPU-bound tiers
			& Cheaper than D for pure compute \\
			\addlinespace
			\textbf{L-series} & Storage optimised
			& High IOPS, local NVMe
			& Big Data, NoSQL, warehousing
			& Maximises disk, not CPU/RAM \\
			\addlinespace
			\textbf{N-series} & GPU accelerated
			& Specialised GPU
			& AI training/inference, rendering
			& Highest tier; GPU choice critical \\
			\bottomrule
		\end{tabular}
	\end{frame}
	\begin{frame}[shrink=15]{B-series in focus: the credit model and a decision path}
		\begin{columns}[T]
			\begin{column}{0.48\textwidth}
				\vspace{0.2cm}
				\textbf{\large How burstable billing works}
				
				\vspace{0.3cm}
				\begin{itemize}
					\setlength{\itemsep}{0.8em}
					\item B-series uses \textbf{CPU credits} for performance.
					\item Below baseline $\rightarrow$ \textbf{accumulate} credits.
					\item Need more power $\rightarrow$ \textbf{burst} up to 100\% vCPU.
					\item Credits exhausted $\rightarrow$ \textbf{throttled} to baseline until they recover.
				\end{itemize}
				
				\vspace{0.4cm}
				\begin{exampleblock}{When to avoid B-series}
					Sustained high CPU (e.g.\ 80\% for hours) exhausts credits and throttles. Use \textbf{D} or \textbf{F} for steady load.
				\end{exampleblock}
			\end{column}
			
			\begin{column}{0.48\textwidth}
				\vspace{0.2cm}
				\textbf{\large A simple decision path}
				
				\vspace{0.3cm}
				\begin{itemize}
					\setlength{\itemsep}{0.8em}
					\item \textbf{Spiky/idle CPU?} \\
					$\hookrightarrow$ Choose \textbf{B-series} (dev/test)
					
					\item \textbf{Steady general workload?} \\
					$\hookrightarrow$ Choose \textbf{D-series} (web/app)
					
					\item \textbf{Needs massive RAM?} \\
					$\hookrightarrow$ Choose \textbf{E-series} (in-memory DB)
					
					\item \textbf{CPU-heavy, low RAM?} \\
					$\hookrightarrow$ Choose \textbf{F-series} (batch, CI)
					
					\item \textbf{Needs high disk IOPS?} \\
					$\hookrightarrow$ Choose \textbf{L-series} (Big Data)
					
					\item \textbf{Needs GPU?} \\
					$\hookrightarrow$ Choose \textbf{N-series} (AI)
				\end{itemize}
			\end{column}
			
		\end{columns}
	\end{frame}
	\begin{frame}{Practical demonstration}
		\begin{itemize}
			\item To create a new Azure Virtual Machine, enter \textit{virtual machines} in search bar.
			\item Under \textbf{Services}, select \textbf{Virtual Machines}.
			\item In the \textbf{Virtual machines} page, select \textbf{Create} from the menu at the top of the page and then select \textbf{Virtual machine} from the drop-down. The \textbf{Create a virtual machine} page opens.
			\item Under \textbf{Instance details}, enter \textit{myVM} for the Virtual machine name and choose \textit{Windows Server 2022 Datacenter: Azure Edition - x64 Gen 2} for the Image. The page will fill in the other fields with defaults and use the VM name to create similar names for the other resources that are needed.
		\end{itemize}
	\end{frame}
	\begin{frame}{Practical demonstration}
		\includegraphics[width=\textwidth]{img/1.png}
	\end{frame}
	\begin{frame}{Practical demonstration}
		\begin{itemize}
			\item Under \textbf{Administrator account}, provide a username, such as \textit{azureuser} and a password. The password must be at least 12 characters long and meet the defined complexity requirements.
		\end{itemize}
		\includegraphics[width=\textwidth]{img/2.png}
	\end{frame}
	\begin{frame}{Practical demonstration}
		\begin{itemize}
			\item Under \textbf{Inbound port rules}, choose \textbf{Allow selected ports} and then select \textbf{RDP (3389)} and \textbf{HTTP (80)} from the drop-down.
		\end{itemize}
		\includegraphics[width=\textwidth]{img/3.png}
	\end{frame}
	\begin{frame}{Practical demonstration}
		\begin{itemize}
			\item At the top of the page, select the \textbf{Management} tab.
			\item In the \textbf{Auto-shutdown} section, select the \textbf{Enable auto-shutdown} option.
			\item Select a \textbf{Shutdown time} and the \textbf{Time zone} to use (UTC is the default). Optionally you can choose to be notified by e-mail before the shutdown.
		\end{itemize}
	\end{frame}
	\begin{frame}{Practical demonstration}
		\begin{itemize}
			\item Leave the remaining defaults and then select the \textbf{Review + create} button at the bottom of the page.
		\end{itemize}
		\includegraphics[width=\textwidth]{img/4.png}
	\end{frame}
	\begin{frame}{Practical demonstration}
		\begin{itemize}
			\item After validation runs, select the \textbf{Create} button at the bottom of the page.
		\end{itemize}
		\includegraphics[width=\textwidth]{img/5.png}
	\end{frame}
	\begin{frame}{Practical demonstration}
		\begin{itemize}
			\item After deployment is complete, select \textbf{Go to resource}.
		\end{itemize}
		\includegraphics[width=\textwidth]{img/6.png}
	\end{frame}
	\begin{frame}{Practical demonstration}
		Connect to virtual machine
		\begin{itemize}
			\item On the overview page for your virtual machine, select \textbf{Connect} and then select \textbf{Connect} again from the drop-down.
			\item In the \textbf{Native RDP} page, select the \textbf{Check access} button to make sure port 3389 is available.
			\item Select \textbf{Download RDP file} to download the connection file to your computer.
			\item Open the downloaded RDP file and click \textbf{Connect} when prompted.
			\item You may receive a certificate warning during the sign-in process. Click \textbf{Yes} or \textbf{Continue} to create the connection.
		\end{itemize}
	\end{frame}
	\begin{frame}[shrink=5]{Screenshots and observations}
		\begin{columns}[T]
			\begin{column}{0.52\textwidth}
				\textbf{\large What the deployment produced}
				\begin{itemize}
					\item \textbf{One form, six resources.} With the name \textit{myVM} the portal also created a NIC, public IP, network security group, virtual network and OS disk -- a VM is never an isolated object.
					\item \textbf{About one minute to a running server}, including validation. The same on physical hardware means ordering, racking, cabling and an OS installation.
					\item \textbf{The size chose itself:} \texttt{Standard\_B2s} (2 vCPU, 4\,GB). Burstable is right for a lab and wrong for a 24/7 database.
				\end{itemize}
			\end{column}
			
			\begin{column}{0.45\textwidth}
				\textbf{\large What we took away}
				\begin{itemize}
					\item \textbf{Only two ports were exposed} (RDP 3389, HTTP 80). RDP open to the Internet is acceptable for a lab, not for production -- Azure Bastion or just-in-time access should replace it.
					\item \textbf{Auto-shutdown at 19:00 UTC} with an e-mail warning deallocates the VM overnight and stops compute billing.
					\item \textbf{Billing follows the VM state, not the workload:} a machine stopped inside Windows still bills; only \textit{Stop (deallocate)} ends compute charges.
				\end{itemize}
			\end{column}
		\end{columns}
		
		\vspace{0.25cm}
		\begin{block}{Observation used later in this deck}
			The VM was cheap only because it was small \emph{and} switched off for most of the day. That single fact -- control over uptime -- decides most of the Azure-versus-own-hardware comparison.
		\end{block}
	\end{frame}
	\section{Azure or own hardware?}
	\begin{frame}[shrink=5]{When buying capacity beats owning hardware}
		\centering
		\setlength{\fboxsep}{5pt}
		\setlength{\fboxrule}{0.6pt}
		
		\colorbox{CyberBurgundy}{\color{white}\small\textbf{One workload \quad$\rightarrow$\quad three questions}}
		
		\vspace{0.3cm}
		
		\begin{tabular}{@{}c@{\hspace{0.15cm}}c@{\hspace{0.15cm}}c@{}}
			\textcolor{CyberRed}{\Large$\downarrow$} &
			\textcolor{CyberRed}{\Large$\downarrow$} &
			\textcolor{CyberRed}{\Large$\downarrow$} \\[0.15cm]
			\fcolorbox{CyberBurgundy}{CyberLight}{\begin{minipage}[t]{3.4cm}
					\centering\footnotesize
					\textcolor{CyberBurgundy}{\textbf{Is the load variable\\ or unknown?}}\\[3pt]
					Proof of concept, season peaks, short projects, demand that may disappear
					\\[5pt]\textcolor{CyberRed}{\textbf{$\rightarrow$ rent a VM}}
			\end{minipage}} &
			\fcolorbox{CyberBurgundy}{CyberLight}{\begin{minipage}[t]{3.4cm}
					\centering\footnotesize
					\textcolor{CyberBurgundy}{\textbf{Is it steady, high and\\ predictable for years?}}\\[3pt]
					Same machine, 24/7, utilisation above roughly half of its capacity
					\\[5pt]\textcolor{CyberRed}{\textbf{$\rightarrow$ own the hardware}}
			\end{minipage}} &
			\fcolorbox{CyberBurgundy}{CyberLight}{\begin{minipage}[t]{3.4cm}
					\centering\footnotesize
					\textcolor{CyberBurgundy}{\textbf{Must the data or latency\\ stay on site?}}\\[3pt]
					Classified data, air-gapped network, sub-millisecond links to plant systems
					\\[5pt]\textcolor{CyberRed}{\textbf{$\rightarrow$ own or dedicated}}
			\end{minipage}}
		\end{tabular}
		
		\vspace{0.35cm}
		\begin{block}{Rule of thumb}
			Rent by default; own only what is measured, steady and long-lived. In this project the workload was neither, so a pay-as-you-go VM with auto-shutdown was the correct answer.
		\end{block}
	\end{frame}
	\begin{frame}[shrink=15]{Total cost of ownership, not sticker price}
		\centering
		\renewcommand{\arraystretch}{1.2}
		\setlength{\tabcolsep}{4pt}
		\footnotesize
		
		\begin{tabular}{@{} L{4.9cm} L{2.9cm} L{3.2cm} @{}}
			\toprule
			\textbf{Cost element (per month, EUR)} & \textbf{Own server} & \textbf{Azure VM, pay-as-you-go} \\
			\midrule
			Upfront purchase & 2\,500--3\,500 (once) & 0 \\
			Hardware amortisation, 4 years & 50--70 & in compute price \\
			Compute: 2 vCPU, 4--8\,GB & -- & 30--70 \\
			Power, cooling, network, IP & 75--115 & in compute price \\
			Backup, spare parts, UPS & 15--30 & 5--20 (storage) \\
			Administration, 2--4\,h/month & 80--240 & 20--50 \\
			\midrule
			\textbf{Indicative monthly total} & \textbf{220--470} & \textbf{75--180} \\
			\textbf{Indicative 3-year total} & \textbf{8\,000--17\,000} & \textbf{2\,700--6\,500} \\
			\bottomrule
		\end{tabular}
		
		\vspace{0.15cm}
		{\scriptsize West Europe, Linux/Windows mix, VAT excluded (Windows adds 15--30 on Azure); re-check with the Azure Pricing Calculator.}
		
		\vspace{0.15cm}
		\begin{block}{Verdict}
			{\footnotesize
				At this size Azure wins on money and on risk: no capex, no admin hours, and auto-shutdown removes the idle cost an owned server pays anyway. Reservations (30--40\% off for one year, up to 60\% for three), spot capacity and Hybrid Benefit close most of the remaining gap. \textbf{Owning wins only for large, steady load} -- 8--16 vCPU or more, 24/7, for three years or longer -- with the room, power and staff already in place.}
		\end{block}
	\end{frame}
	\section{Results}
	\begin{frame}{Key findings}
		\begin{itemize}
			\item \textbf{Flexibility and control}
			\begin{itemize}
				\item Azure VMs provide full control over the operating system,
				installed software and runtime environment.
				\item They support different operating systems, architectures
				and workload requirements.
			\end{itemize}
			
			\item \textbf{Scalability and availability}
			\begin{itemize}
				\item VM sizes can be selected according to CPU, memory and
				storage requirements.
				\item Resources can be scaled up or down as workload demands change.
			\end{itemize}
			
			\item \textbf{Cost optimisation}
			\begin{itemize}
				\item Different VM families allow optimisation for specific workloads.
				\item B-series VMs can reduce costs for workloads with variable CPU usage.
				\item Auto-shutdown helps avoid unnecessary compute costs.
			\end{itemize}
			
			\item \textbf{Cloud integration}
			\begin{itemize}
				\item VMs can be integrated with other Azure services,
				networks and storage solutions.
			\end{itemize}
		\end{itemize}
	\end{frame}
	
	\begin{frame}{Practical advantages and limitations}
		\begin{columns}[T]
			\begin{column}{0.48\textwidth}
				\textbf{\large Advantages}
				\begin{itemize}
					\item No need to purchase or maintain physical hardware.
					\item Full control over the operating system.
					\item Support for custom software and configurations.
					\item Flexible resource selection.
					\item Integration with the Azure ecosystem.
					\item Ability to automate deployment and management.
				\end{itemize}
			\end{column}
			
			\begin{column}{0.48\textwidth}
				\textbf{\large Limitations}
				\begin{itemize}
					\item Operating system updates and security remain the user's responsibility.
					\item Costs can increase if resources are not managed properly.
					\item Incorrect VM sizing can cause poor performance or overspending.
					\item Network and access configuration require attention.
					\item B-series performance depends on available CPU credits.
				\end{itemize}
			\end{column}
		\end{columns}
	\end{frame}
	
	
	\begin{frame}{Lessons learned}
		\begin{itemize}
			\item \textbf{Choosing the right VM size matters}
			\begin{itemize}
				\item VM selection should be based on actual workload requirements,
				not simply the highest available specifications.
			\end{itemize}
			
			\item \textbf{Security must be considered from the beginning}
			\begin{itemize}
				\item Only necessary inbound ports should be exposed.
				\item Strong authentication and regular security updates are essential.
			\end{itemize}
			
			\item \textbf{Cost management is important}
			\begin{itemize}
				\item Auto-shutdown can prevent unnecessary spending.
				\item VM deallocation should be considered when resources are no longer needed.
			\end{itemize}
			
			\item \textbf{IaaS requires more administration}
			\begin{itemize}
				\item Azure manages the underlying infrastructure,
				but users remain responsible for the guest operating system.
			\end{itemize}
		\end{itemize}
		
		\vspace{0.3cm}
		\begin{block}{Main takeaway}
			Azure Virtual Machines offer flexibility and control,
			but efficient usage requires proper sizing, security
			configuration and continuous cost management.
		\end{block}
	\end{frame}
	
	\section{Conclusions and sources}
	\begin{frame}{Conclusions: Key results and Azure value}
		\begin{columns}[T]
			\begin{column}{0.48\textwidth}
				\textbf{\large Key project results}
				\begin{itemize}
					\setlength{\itemsep}{0.5em}
					\item \textbf{Rapid provisioning:} One portal form deployed six linked resources (VM, NIC, public IP, NSG, VNet, OS disk) in about one minute.
					\item \textbf{Full OS control:} IaaS delivered the flexibility required for custom runtimes and software that does not fit PaaS.
					\item \textbf{Lower TCO for variable load:} Pay-as-you-go with auto-shutdown (75--180\,EUR/month) beat an owned server (220--470\,EUR/month) with zero upfront capex.
				\end{itemize}
			\end{column}
			
			\begin{column}{0.48\textwidth}
				\textbf{\large How Azure services supported the work}
				\begin{itemize}
					\setlength{\itemsep}{0.5em}
					\item \textbf{Burstable B-series:} Matched an idle-heavy lab workload using CPU credits without paying for unused capacity.
					\item \textbf{Built-in cost control:} Scheduled auto-shutdown deallocated the VM overnight, directly cutting compute charges.
					\item \textbf{Integrated networking:} Network Security Groups enabled quick port filtering (RDP 3389, HTTP 80) during setup.
				\end{itemize}
			\end{column}
		\end{columns}
		
		\vspace{0.25cm}
		\begin{block}{Core verdict}
			For short-lived, variable, or proof-of-concept workloads, renting an Azure VM is faster, cheaper, and lower-risk than purchasing and maintaining physical hardware.
		\end{block}
	\end{frame}
	
	\begin{frame}{Conclusions: Limitations and next steps}
		\begin{columns}[T]
			\begin{column}{0.48\textwidth}
				\textbf{\large Practical limitations and lessons}
				\begin{itemize}
					\setlength{\itemsep}{0.5em}
					\item \textbf{Shared responsibility:} Azure handles hardware, but OS patching, updates, and hardening remain the user's job.
					\item \textbf{Deallocation matters:} Shutting down inside the OS still incurs compute billing; only \textit{Stop (deallocate)} stops charges.
					\item \textbf{Sizing trade-offs:} B-series throttles once CPU credits run out, making D- or F-series necessary for sustained loads.
					\item \textbf{Exposure risk:} Leaving RDP open to the public Internet is acceptable only for isolated testing.
				\end{itemize}
			\end{column}
			
			\begin{column}{0.48\textwidth}
				\textbf{\large Recommended next steps}
				\begin{itemize}
					\setlength{\itemsep}{0.5em}
					\item \textbf{Harden access:} Replace public RDP with Azure Bastion, VPN, or just-in-time VM access.
					\item \textbf{Automate deployments:} Export ARM/Bicep templates to recreate environments consistently and without manual errors.
					\item \textbf{Add observability:} Enable Azure Monitor to track VM health, CPU credit balance, and network dependencies.
				\end{itemize}
			\end{column}
		\end{columns}
	\end{frame}
	
	\begin{frame}[shrink=10]{References and resources}
		\small
		\begin{thebibliography}{10}
			\bibitem{azure} Microsoft Learn, \emph{Microsoft Azure documentation}, \url{https://learn.microsoft.com/en-us/azure/}.
			\bibitem{ai} Microsoft Learn, \emph{Azure AI services documentation}, \url{https://learn.microsoft.com/en-us/azure/ai-services/}.
			\bibitem{paas} Microsoft \url{https://azure.microsoft.com/en-us/resources/cloud-computing-dictionary/what-is-paas}
			\bibitem{iaas} Microsoft \url{https://azure.microsoft.com/en-us/resources/cloud-computing-dictionary/what-is-iaas/}
			\bibitem{vmsize} Microsoft \url{https://learn.microsoft.com/en-us/azure/virtual-machines/sizes/overview}
			\bibitem{bcpu} \url{https://learn.microsoft.com/en-us/azure/virtual-machines/sizes/b-series-cpu-credit-model}
			\bibitem{pricing} Azure pricing page and calculator, \url{https://azure.microsoft.com/en-us/pricing/calculator/}
			\bibitem{reserved} Microsoft Learn, \emph{Azure Reservations and savings plans}, \url{https://learn.microsoft.com/en-us/azure/cost-management-billing/reservations/}
			\bibitem{hybrid} Microsoft Learn, \emph{Azure Hybrid Benefit}, \url{https://learn.microsoft.com/en-us/azure/azure-hybrid-benefit/}
		\end{thebibliography}
	\end{frame}
	\begin{frame}{Questions and discussion}
		\centering
		{\LARGE\bfseries Questions?}
		\vspace{0.6cm}
		
		{\large We welcome your questions and feedback.}
	\end{frame}
	\begin{frame}[plain]
		\centering
		{\Huge\bfseries Thank you!}
	\end{frame}
\end{document}